\section{Alignment with Standards and Regulations}
	
	Security and privacy requirements do not exist in isolation: they are guided by community best practices, technical standards and legal frameworks.
	This section places ViperVault in relation to the most relevant reference points, namely the OWASP guidelines, the NIST recommendations and European regulations such as GDPR and NIS2.
	By aligning the threat model with these standards, the project demonstrates not only technical robustness but also regulatory compliance and long-term sustainability.
	
	\subsection{OWASP and NIST}
		
		ViperVault has been designed with explicit reference to the \textbf{OWASP Mobile Application Security Verification Standard (MASVS)} and the \textbf{Mobile Security Testing Guide (MSTG)}.
		Core requirements such as secure credential storage (MASVS-STORAGE-2), strong and up-to-date cryptography, secure input handling and the wiping of sensitive information from memory are explicitly addressed.
		This ensures that the application does not merely rely on cryptographic primitives but also respects best practices for their implementation and for the surrounding environment.
		In addition, the application complies with recommendations from \textbf{NIST SP 800-63B} on digital authentication.
		Specifically, the use of \textbf{Argon2id} as a memory-hard key derivation function satisfies the requirement to resist large-scale GPU or ASIC cracking attempts, while the password policy enforces sufficient entropy, prohibits commonly used secrets and mitigates risks associated with password reuse.
		Together, these controls position ViperVault in line with modern expectations for secure identity and credential management.
	
	\subsection{GDPR}
		
		From a legal and privacy perspective, the design of ViperVault is strongly aligned with the \textbf{General Data Protection Regulation (GDPR)}.
		The application follows the principle of \textbf{Privacy by Design}, meaning that it collects and processes only the minimum data required for its function.
		All vault contents are encrypted client-side, no personal identifiers are stored in plaintext and no unnecessary telemetry or logs are generated.
		The \textbf{right to erasure} is also inherently respected: once a vault is deleted, it is cryptographically unrecoverable, ensuring that the user retains full control over their data lifecycle.
		Furthermore, the principle of \textbf{data minimization} is applied rigorously, with no additional metadata (such as usage statistics or access logs) being collected.
		This not only protects users from privacy risks but also simplifies compliance with regulatory audits.
	
	\subsection{NIS2}
		
		Furthermore, ViperVault is compliant with the European \textbf{NIS2 Directive}, which expands cybersecurity obligations for digital service providers.
		Even though the project is currently developed by a single developer, mechanisms are planned to support a security management lifecycle.
		These include \textbf{reproducible builds} to ensure the integrity of released binaries, a process for responsible \textbf{vulnerability disclosure} and the use of \textbf{digitally signed updates} to prevent tampering.
		In terms of incident response, the project prioritizes transparency and agility.
		Public transparency logs ensure that update integrity can be independently verified and a lightweight process for rapid patching helps minimize exposure when new vulnerabilities are discovered.
		In this way, even with limited resources, ViperVault can embody the spirit of NIS2: resilience, accountability and continuous improvement in security.